\begin{textsample}{POS Dim 2 – persona\_gemini – Score 20.00 – t571\_gemini.txt}  \label{ex:f2_pos_045}
The \textbf{planet} is sending a clear signal.
June was the hottest month ever recorded, and July is on track to follow suit.
We are seeing the consequences in real-time, with extreme \textbf{heatwaves} scorching \textbf{communities} across the globe.
In Europe, temperatures in France have soared to 45.6°C, while in India, temperatures exceeding 50°C have caused deaths and raised the terrifying \textbf{possibility} that parts of the country could become uninhabitable in the \textbf{future}.
By 2050, 75 % of the \textbf{world} 's cities are projected to \textbf{face} shifted \textbf{climates} as heat exposure continues to rise.
A primary \textbf{driver} of this \textbf{crisis} is \textbf{coal}, which is responsible for 46 % of global CO2 emissions.
Despite this, European \textbf{leaders} are \textbf{delaying} the essential transition away from fossil \textbf{fuels}, citing economic concerns.
This hesitation ignores the \textbf{reality} that a shift to net zero emissions could create jobs, yield massive health benefits, and save trillions of dollars.
But \textbf{people} are not waiting for politicians to \textbf{act}.
Global protests are \textbf{demanding} change, from the successful movement to save Germany 's Hambach Forest to the worldwide Fridays for \textbf{Future} strikes.
To help connect this growing opposition, Greenpeace 's Unite for \textbf{Climate} tour is travelling across Europe to expose the \textbf{barriers} to progress and highlight the devastating \textbf{impacts} of the \textbf{climate} \textbf{crisis} on \textbf{communities}.

% matched lemmas: act, barrier, climate, coal, community, crisis, delay, demand, driver, face, fuel, future, heatwave, impact, leader, people, planet, possibility, reality, world
\end{textsample}
